\chapter{Sprint 3: Web Dashboard and Live Monitoring}
\chaptertoc

\clearpage

\section{Introduction}
Sprint 3 covered the frontend work from the first design decisions and page structure to the live backend integration that made the dashboard operational. The sprint started by defining the visual direction and navigation model, then connected the interface to backend services and runtime events so the dashboard could react to live data, transport changes, and feed lifecycle actions.

This two-phase approach was intentional. Settling the design system and page skeleton first meant that the live integration work could focus on behavior and reliability instead of redesigning the interface at the same time.

\section{Sprint Objectives}
This sprint was organized around six outcomes spanning design and live integration:

\textbf{Design and scaffolding objectives:}
\begin{itemize}
    \item define a coherent visual language for the dashboard,
    \item prepare the main page layouts before live data integration,
    \item build reusable UI structures that can support later features.
\end{itemize}

\textbf{Live integration objectives:}
\begin{itemize}
    \item make the dashboard consume the backend contract without ad hoc transformations,
    \item keep the interface responsive when live transport is available, delayed, or interrupted,
    \item validate that monitoring, feed management, and fallback behavior remain understandable for the operator.
\end{itemize}

\section{Sprint Backlog}
The Sprint 3 backlog separates manager-facing user stories from technical implementation tasks. User stories represent stakeholder needs and value; technical tasks are internal infrastructure concerns.

\subsection{User Stories}
The user stories in Sprint 3 represent manager and admin needs for dashboard interfaces and live monitoring, as shown in Table~\ref{tab:sprint3_user_stories}.

\begin{table}[!htbp]
    \centering
    \footnotesize
    \setlength{\tabcolsep}{3pt}
    \renewcommand{\arraystretch}{0.88}
    \begin{tabularx}{\textwidth}{|l|X|l|X|}
        \hline
        User StoryID & User Story & TaskID & Task \\ \hline
        US3.1 & As a manager, I want a consistent visual style so that the dashboard feels like one coherent product. & T3.1 & Define colors, typography, spacing, and reusable component rules. \\ \hline
        US3.2 & As a manager, I want the main page layouts prepared so that onboarding, monitoring, editing, and analytics screens are available. & T3.2 & Prepare the landing page, dashboard, setup wizard, zone editor, analytics, and settings layouts. \\ \hline
        US3.3 & As an admin, I want to create and maintain manager accounts so that system access is centrally controlled. & T3.3 & Build the admin dashboard for manager account creation, editing, and deletion. \\ \hline
        US3.4 & As a manager, I want the dashboard connected to backend responses so that the interface reflects live operational data. & T3.4 & Implement the API integration layer and typed request handling. \\ \hline
        US3.5 & As a manager, I want to manage feeds from the interface so that feed lifecycle actions stay centralized. & T3.5 & Connect login state, protected routes, and feed control actions to the dashboard. \\ \hline
        US3.6 & As a manager, I want live updates without refreshing so that changes are visible immediately. & T3.6 & Add websocket-driven live state updates and dashboard orchestration. \\ \hline
        US3.7 & As a manager, I want a suitable video transport chosen automatically so that the feed remains usable when conditions change. & T3.7 & Add transport selection, fallback streaming, and reconnection logic. \\ \hline
        US3.8 & As a manager, I want the dashboard to stay usable during interruptions so that monitoring does not stop during transient failures. & T3.8 & Add fallback polling and stale-data guards. \\ \hline
    \end{tabularx}
    \caption{Sprint 3 user stories}
    \label{tab:sprint3_user_stories}
\end{table}
\normalsize

\subsection{Technical Tasks}
The technical tasks in Sprint 3 represent internal implementation concerns needed to support the user stories.

\begin{table}[!htbp]
    \centering
    \footnotesize
    \setlength{\tabcolsep}{3pt}
    \renewcommand{\arraystretch}{0.88}
    \begin{tabularx}{\textwidth}{|l|X|l|X|}
        \hline
        Task ID & Task & Subtask ID & Subtask \\ \hline
        T3.10 & Set up the frontend project with shared components and page scaffolding so that the interface can grow without duplication. & & Create reusable UI component library and establish project structure. \\ \hline
        T3.11 & Validate dashboard setup and onboarding flows with integration tests to ensure the integration behaves correctly end to end. & & Implement scenario-based tests for dashboard lifecycle and error paths. \\ \hline
    \end{tabularx}
    \caption{Sprint 3 technical tasks}
    \label{tab:sprint3_technical_tasks}
\end{table}
\normalsize

\section{Design}

\subsection{UI and UX Design}
The project design specification introduces a ``Terminal Minimal'' style and a seven-page navigation model. The goal was to give the dashboard a simple visual identity that fits a monitoring application. The design language emphasizes:
\begin{itemize}
    \item dark operational surfaces optimized for continuous monitoring,
    \item strong semantic colors for status and alert levels,
    \item card-based information density with consistent spacing and hierarchy,
    \item explicit workflow guidance in setup and zone-editing screens.
\end{itemize}

The designed page set includes the landing page, dashboard, setup wizard, zone editor, analytics, settings, and heatmap views. These screens helped define what the frontend had to support before live data arrived.

\subsection{Sequence Diagrams}
The Sprint 3 sequence for dashboard updates and alert surfacing is illustrated in Figure~\ref{fig:sprint3_sequence}.
\begin{figure}[H]
    \centering
    \includegraphics[width=0.95\textwidth]{\detokenize{chapter5_sequence.pdf}}
    \caption{Sprint 3 sequence diagram for dashboard updates and alert surfacing.}
    \label{fig:sprint3_sequence}
\end{figure}

Three main flows were implemented:
\begin{enumerate}
    \item \textbf{Feed lifecycle:} a dashboard action updates the feed state and the new status is sent back to the interface.
    \item \textbf{Live metrics:} metric updates are received from the backend and reflected immediately on the dashboard.
    \item \textbf{Alert surfacing:} alerts reach the interface and appear in the activity and attention areas.
\end{enumerate}

\subsection{Dashboard Data-Flow Sequence Diagram}
Figure~\ref{fig:sprint3_dashboard_data_flow_sequence} details how data moves between the user, the React dashboard, the API layer, the backend, the runtime worker, the websocket channel, and the media transport layer.
\begin{figure}[H]
    \centering
    \includegraphics[width=0.95\textwidth]{\detokenize{chapter5_dashboard_data_flow_sequence.pdf}}
    \caption{Sprint 3 detailed sequence diagram for dashboard data flow, live events, fallback polling, and video transport.}
    \label{fig:sprint3_dashboard_data_flow_sequence}
\end{figure}

This sequence starts with protected dashboard access and initial REST queries for feeds and system health. The dashboard then opens the metrics websocket, receives an initial snapshot, and keeps its local state synchronized from feed lifecycle, metric, alert, and warning events. Feed actions are sent as REST commands and reflected back through websocket status events. When live events become stale, the frontend polls the feed list as a fallback. For video playback, the dashboard negotiates WebRTC through the backend when the transport is ready and otherwise uses the MJPEG stream endpoint as the simpler fallback.

\subsection{Live Dashboard Activity Diagram}
The frontend runtime behavior is summarized in Figure~\ref{fig:sprint3_live_activity}. This diagram complements the sequence view by showing how the dashboard decides between real-time websocket updates, fallback polling, feed actions, and video transport modes.
\begin{figure}[H]
    \centering
    \includegraphics[width=0.78\textwidth]{\detokenize{chapter5_live_dashboard_activity.pdf}}
    \caption{Sprint 3 live dashboard activity diagram for websocket updates, fallback polling, feed actions, and transport selection.}
    \label{fig:sprint3_live_activity}
\end{figure}

The activity begins with authenticated dashboard access, then loads the feed list and system health. After that, the websocket channel keeps feed status, metric updates, alerts, and warnings synchronized with the local query cache. If the websocket becomes stale or unavailable while feeds are running, the dashboard temporarily polls the feed endpoint so the monitoring surface remains usable. In parallel, each feed card selects the best available video surface: WebRTC when the backend reports a ready stream, MJPEG when WebRTC is unavailable or fails, and a stored preview when the feed is not yet running.

\subsection{Design Rationale}
The design phase was kept as a foundation step rather than a full UI release. The goal was to settle navigation, component boundaries, and interaction patterns before live data and transport logic arrived. A clean shell was a useful deliverable because the integration work depended on it.

\section{Implementation}

\subsection{Project Setup and Layout Skeleton}
The implementation started by turning the design specification into an application skeleton. The workspace setup established the application entry point, protected navigation, the dashboard layout, and the shared components that all pages reuse.

The order of the pages was not arbitrary. The layout had to support a workflow that begins with onboarding, continues through configuration, and ends with live supervision. That is why the page skeleton was treated as an operational path rather than a list of unrelated screens.

\subsection{Styling and Reusable Interface Rules}
The visual system was defined once and then applied consistently across the dashboard. Spacing, typography, surface treatment, and semantic status colors were treated as reusable rules rather than per-page decisions. This reduced design drift: when a status card, a table, or a setup control appears in several places, it behaves like the same product element instead of being restyled each time.

\subsection{API Integration}
The frontend relies on a dedicated API layer that keeps the communication with the backend structured. It defines the main data types used by the dashboard and centralizes request handling so errors can be managed consistently. This layer acts as the contract boundary: it prevents the UI from depending directly on backend implementation details and makes it easier to reason about feed lists, dashboard summaries, and management actions as explicit typed responses.

\subsection{Authentication and Role-Based Navigation}
Sprint 3 connected the login and sign-up flows to the dashboard. The auth provider clears cached data when the session changes, protected routes redirect users according to their role, and administrators are taken to the manager administration page while managers are routed to the monitoring dashboard. This keeps the frontend aligned with the backend access rules and avoids exposing the wrong interface to the wrong user, with the manager sign-in screen shown in Figure~\ref{fig:sprint3_login_page}.
\begin{figure}[H]
    \centering
    \includegraphics[width=0.65\textwidth]{\detokenize{LoginPage.png}}
    \caption{Manager sign-in screen used to access the dashboard in Sprint 3.}
    \label{fig:sprint3_login_page}
\end{figure}
The complementary manager account creation screen introduced in the same authentication flow is presented in Figure~\ref{fig:sprint3_signup_page}.
\begin{figure}[H]
    \centering
    \includegraphics[width=0.65\textwidth]{\detokenize{SignupPage.png}}
    \caption{Manager account creation screen introduced as part of the Sprint 3 authentication flow.}
    \label{fig:sprint3_signup_page}
\end{figure}

\subsection{Dashboard State Orchestration}
A dedicated dashboard state layer coordinates data fetching and updates. From the received feed data, the dashboard derives indicators such as online feeds, total people, average wait time, and attention counts, then exposes them to the widgets.

The state layer also helps absorb transient data gaps. When a websocket event arrives late, the dashboard can keep displaying the last valid snapshot instead of dropping the screen into an empty state. That design choice improves usability during unstable network periods and keeps the monitoring surface trustworthy.

\subsection{Live Update Reliability}
The live-update layer includes:
\begin{itemize}
    \item reconnection delays,
    \item a heartbeat check,
    \item a guard against stale data,
    \item fallback polling when live updates fail,
    \item grouped updates to keep the interface responsive.
\end{itemize}
When the websocket connection drops, the dashboard falls back to periodic polling instead of freezing, so operators still have a usable view during transient issues.

This fallback behavior was treated as a first-class feature rather than an error path. In a queue-monitoring context, the operator needs continuity more than perfect transport purity, so the UI must remain informative even when the preferred live path is temporarily unavailable.

\subsection{Live Video Delivery}
Live feed cards use a transport decision layer to choose the most suitable viewing mode. For each running feed, the frontend checks whether live viewing is ready. If the higher-quality live stream is available, it is used first; otherwise the dashboard falls back to the simpler video stream.

The live-streaming flow is handled through the browser's live video connection support \cite{webrtc_w3c}:
\begin{itemize}
    \item create a connection for receiving video,
    \item gather the network information needed for signaling,
    \item send the browser offer to the backend,
    \item apply the backend response and switch the feed card to live state.
\end{itemize}

Reconnection uses increasing delays between attempts. When the problem is not recoverable, the system stops retrying and switches to the fallback stream so the dashboard remains usable.

When server-side annotations are active, the dashboard hides local overlays to avoid drawing the same boxes twice on the same video.

\subsection{Dashboard UI Integration}
The dashboard page brings together live feed controls, setup dialogs, source onboarding, KPI visualizations, and charts. The onboarding workflow supports both uploaded videos and ONVIF camera sources. For camera sources, the manager discovers devices, applies per-camera credentials when needed, resolves RTSP streams, tests the selected camera, captures a snapshot for zone tracing, and then stages the feed for batch launch. The emphasis was on making those elements work together as one operational interface.

The practical result is that a manager can move from onboarding to monitoring without leaving the dashboard context. That continuity mirrors the workflow described in the earlier requirements chapter, while Figure~\ref{fig:sprint3_admin_dashboard} shows the administrator screen used for creating and maintaining manager accounts.
\begin{figure}[H]
    \centering
    \includegraphics[width=0.9\textwidth]{\detokenize{AdminDashboarUi.png}}
    \caption{Administrator manager-administration screen for creating and maintaining manager accounts.}
    \label{fig:sprint3_admin_dashboard}
\end{figure}
The manager-facing monitoring dashboard, including the live surface, feed summary cards, and recent activity panel, is presented in Figure~\ref{fig:sprint3_manager_dashboard}.
\begin{figure}[H]
    \centering
    \includegraphics[width=0.95\textwidth]{\detokenize{ManagerDashboardUi.png}}
    \caption{Manager dashboard showing the live monitoring surface, feed summary cards, and recent activity panel.}
    \label{fig:sprint3_manager_dashboard}
\end{figure}
The onboarding workflow starts with the feed setup screen shown in Figure~\ref{fig:sprint3_feed_config}, where the manager prepares an uploaded video or an ONVIF camera source before zone tracing and model selection.
\begin{figure}[H]
    \centering
    \includegraphics[width=0.9\textwidth]{\detokenize{FeedConfigUi.png}}
    \caption{Feed setup screen used to prepare an uploaded video or ONVIF camera source before zone tracing and model selection.}
    \label{fig:sprint3_feed_config}
\end{figure}
The queue polygon for the selected feed is defined through the zone tracing screen shown in Figure~\ref{fig:sprint3_zone_trace}.
\begin{figure}[H]
    \centering
    \includegraphics[width=0.9\textwidth]{\detokenize{ZoneTraceUi.png}}
    \caption{Zone tracing screen used to define the queue polygon for the selected feed.}
    \label{fig:sprint3_zone_trace}
\end{figure}
Before launching the feed, the manager selects the YOLO model size through the screen shown in Figure~\ref{fig:sprint3_model_select}.
\begin{figure}[H]
    \centering
    \includegraphics[width=0.9\textwidth]{\detokenize{ModelSelectUi.png}}
    \caption{Model selection screen used to choose the YOLO size before launching the feed.}
    \label{fig:sprint3_model_select}
\end{figure}
The staged feed or batch of feeds is then confirmed in the review and launch panel presented in Figure~\ref{fig:sprint3_review_panel}.
\begin{figure}[H]
    \centering
    \includegraphics[width=0.95\textwidth]{\detokenize{ReviewPanelUi.png}}
    \caption{Review and launch panel used to confirm staged feeds before batch creation or launch.}
    \label{fig:sprint3_review_panel}
\end{figure}

\subsection{Analytics and Historical Reporting}
The analytics view extends the dashboard from live supervision into historical reporting. It aggregates archived alert history and queue metrics into a summary that helps managers identify trend shifts and recurring bottlenecks without leaving the web application.

Four KPI cards provide a fast overview: average wait time, peak queue length, stability score, and total alerts across the selected range. A period selector (hourly, daily, weekly) controls the aggregation granularity, while date filters allow managers to isolate specific operating windows. The page then visualizes wait time and queue size trends, highlights the top-performing and underperforming zones, and summarizes the alert mix by severity.

The analytics view also supports CSV export so the same statistics can be used in external reporting. When the archive is empty, the interface seeds a mock dataset derived from active feeds to keep the dashboard informative during early pilots. Figure~\ref{fig:sprint3_analytics_page} presents the consolidated analytics dashboard.
\begin{figure}[H]
    \centering
    \includegraphics[width=0.95\textwidth]{\detokenize{AnalyticsPageUi.png}}
    \caption{Analytics dashboard with KPI summaries, trend charts, zone breakdowns, and export controls.}
    \label{fig:sprint3_analytics_page}
\end{figure}

\subsection{Integration Testing}
Frontend integration behavior is validated through scenario-based tests that cover dashboard setup, live event handling, transport fallback, feed routing, and error behavior.

The tests are useful not just for regression detection, but also for documenting what the sprint is supposed to guarantee. They show that the implementation is not finished when the screen renders; it is finished when the system still behaves correctly under reconnects, route changes, and fallback transitions.

\section{Conclusion}
By the end of Sprint 3, the dashboard could react to live state, switch transport modes, and stay usable when the live connection failed. That completed the frontend delivery and prepared the project for alert automation and final evaluation.

\section{Sprint Retrospective}
This sprint confirmed that separating the design foundation from live integration was the right approach. The most useful outcome was the orchestration layer that connects session state, feed actions, live metrics, and transport fallback into one predictable operator experience.

The main tradeoff was complexity at the boundary between frontend and backend. More logic moved into the dashboard, but that logic was necessary to keep the system resilient when transport conditions changed. In the final review, this was considered acceptable because the result is easier to operate and easier to explain in the report.
